\documentclass[11pt]{article}
\usepackage{docstyle}
% Lesson 3 lost-mark points (from the material plus teaching observation), numbered once
\setcounter{lostmark}{0}
\lostmark{tpcopy}{Copying the shape: the \tturns of $f$ become \kw{$x$-intercepts} of $f'$, not \tturns of $f'$.}
\lostmark{sign}{Wrong side of the axis: where $f$ goes down, $f'$ is \kw{below} the $x$-axis.}
\lostmark{degree}{Wrong type of graph: a parabola's $f'$ is a straight line; a cubic's $f'$ is a parabola.}
\lostmark{steep}{The steepest point of $f$ not matched: it gives the \tturn of $f'$.}
\lostmark{roots}{Reverse sketch: the zeros of $f'$ are the \tturns of $f$, not where $f$ crosses the axis.}
\lostmark{marks}{Key $x$-values not lined up with the given graph, or the new graph not labelled.}

\newlength\figunit
\def\showfp{}
\begin{document}
\handouthead{Lesson 3 \textperiodcentered\ Sketching Gradient Functions}{NCEA Level 2 Mathematics \textperiodcentered\ AS 91262 Apply calculus methods in solving problems \textperiodcentered\ Handout}

\section{The one idea}\label{idea}
\begin{keybox}
The \kw{height} of the $f'$ graph is the \kw{steepness} of the $f$ graph.\\
Where $f$ is flat (a \tturn), $f'$ is on the $x$-axis. Where $f$ goes up, $f'$ is above the axis; where $f$ goes down, $f'$ is below.
\end{keybox}
Always draw the new graph \kw{directly under} the given one, using the same $x$-scale, and line up the key $x$-values with a ruler.

\section{Translating features}\label{table}
\begin{tabular}{@{}p{0.44\linewidth}p{0.50\linewidth}@{}}
\toprule
on the graph of $f$ & on the graph of $f'$\\
\midrule
\tturn (maximum or minimum) & \kw{crosses or touches the $x$-axis} ($f'=0$)\\
increasing (going up) & \kw{above} the $x$-axis ($f'>0$)\\
decreasing (going down) & \kw{below} the $x$-axis ($f'<0$)\\
steepest point (\tinflect) & \tturn of $f'$\\
straight line with gradient $m$ & horizontal line at height $m$\\
parabola & straight line\\
cubic & parabola\\
quartic & cubic\\
\bottomrule
\end{tabular}

\smallskip
The power drops by one: that is why a parabola gives a line and a cubic gives a parabola.

\section{Lines and parabolas}\label{parab}
\begin{minipage}[t]{0.50\linewidth}
A straight line has the same gradient everywhere, so $f'$ is a \kw{horizontal line}.
E.g.\ $f(x)=2x+1$ gives $f'(x)=2$ for every $x$.

\smallskip
A parabola has one \tturn, so $f'$ is a \kw{straight line through that $x$-value}.
For the parabola on the right (minimum at $x=1$): down then up, so $f'$ goes from below the axis to above it,
crossing at $x=1$. Here $f(x)=x^{2}-2x-3$ and $f'(x)=2x-2$.
\end{minipage}\hfill
\begin{minipage}[t]{0.46\linewidth}\centering\vspace{0pt}
\setlength\figunit{6.5mm}
% f(x) = x^2 - 2x - 3 = (x-1)^2 - 4: minimum at x = 1.  Caller sets \figunit.
\sketch[xmin=-2,xmax=4,ymin=-5,ymax=6,unit=\figunit,yunit=0.5\figunit,xstep=\figxstep,ystep=\figystep,name={$f(x)$},
  curve={\x*\x-2*\x-3},dmin=-2,dmax=4,
  extra={\draw[dashed, line width=0.6pt] (1,-5) -- (1,6);}]
\\[2mm]
% f'(x) = 2x - 2: crosses the x-axis at x = 1.  Curve drawn only when \showfp is defined.
\sketch[xmin=-2,xmax=4,ymin=-6,ymax=6,unit=\figunit,yunit=0.5\figunit,xstep=\figxstep,ystep=\figystep,name={$f'(x)$},
  curve={2*\x-2},show=\ifdefined\showfp 1\else 0\fi,dmin=-2,dmax=4,
  extra={\draw[dashed, line width=0.6pt] (1,-6) -- (1,6);}]

\end{minipage}

\section{Cubics and quartics}\label{cubic}
\begin{minipage}[t]{0.50\linewidth}
\subsection{Write it in this order}
\begin{enumerate}[leftmargin=7mm, itemsep=1pt]
\item Mark the $x$-value of each \tturn of $f$ on the new axis: these are the \kw{zeros of $f'$}.
\item In each interval: $f$ up $\to$ $f'$ above; $f$ down $\to$ $f'$ below.
\item The steepest point of $f$ gives the \tturn of $f'$.
\item Join smoothly. Check the type: cubic $\to$ parabola.
\end{enumerate}

\begin{example}{Worked example 1}
$f(x)=x^{3}-3x$ has a maximum at $x=-1$ and a minimum at $x=1$.
\begin{steps}
\item Zeros of $f'$ at $x=-1$ and $x=1$.
\item Up, down, up: $f'$ is above, below, above.
\item Steepest downhill at $x=0$: lowest point of $f'$.
\end{steps}
So $f'$ is an up-parabola: $f'(x)=3x^{2}-3$.
\end{example}
A quartic with three \tturns gives a cubic $f'$ with three zeros.
\end{minipage}\hfill
\begin{minipage}[t]{0.46\linewidth}\centering\vspace{2mm}
\setlength\figunit{7.2mm}
% f(x) = x^3 - 3x: maximum at x = -1, minimum at x = 1, steepest downhill at x = 0.
\sketch[xmin=-3,xmax=3,ymin=-6,ymax=6,unit=\figunit,yunit=0.45\figunit,xstep=\figxstep,ystep=\figystep,name={$f(x)$},
  curve={\x*\x*\x-3*\x},dmin=-2.6,dmax=2.6,
  extra={\draw[dashed, line width=0.6pt] (-1,-6) -- (-1,6) (1,-6) -- (1,6);}]
\\[2mm]
% f'(x) = 3x^2 - 3: zeros at x = -1 and 1, minimum -3 at x = 0.  Curve only with \showfp.
\sketch[xmin=-3,xmax=3,ymin=-4,ymax=8,unit=\figunit,yunit=0.45\figunit,xstep=\figxstep,ystep=\figystep,name={$f'(x)$},
  curve={3*\x*\x-3},show=\ifdefined\showfp 1\else 0\fi,dmin=-2.2,dmax=2.2,
  extra={\draw[dashed, line width=0.6pt] (-1,-4) -- (-1,8) (1,-4) -- (1,8);}]

\end{minipage}

\section{From $f'$ back to $f$}\label{reverse}
\begin{minipage}[t]{0.50\linewidth}
\begin{tabular}{@{}p{0.52\linewidth}p{0.40\linewidth}@{}}
\toprule
$f'$ graph & $f$ graph\\
\midrule
crosses from $+$ to $-$ & maximum\\
crosses from $-$ to $+$ & minimum\\
above the axis & increasing\\
below the axis & decreasing\\
\tturn of $f'$ & steepest point\\
\bottomrule
\end{tabular}

\medskip
The $f'$ graph shows gradients only, so the \kw{height of $f$ is unknown}: sketch \emph{a possible} graph.
Any vertical shift is correct.

\begin{example}{Worked example 2}
$f'$ (top right) crosses from $-$ to $+$ at $x=-1$ and from $+$ to $-$ at $x=3$.
So $f$ has a \kw{minimum at $x=-1$} and a \kw{maximum at $x=3$}; it is a negative cubic.
One possible $f$ is drawn under it.
\end{example}
\end{minipage}\hfill
\begin{minipage}[t]{0.46\linewidth}\centering\vspace{0pt}
\setlength\figunit{7.5mm}
% given f'(x) = -(x+1)(x-3) = -x^2 + 2x + 3: zeros at -1 and 3.
\sketch[xmin=-2,xmax=4,ymin=-6,ymax=5,unit=\figunit,yunit=0.45\figunit,xstep=\figxstep,ystep=\figystep,name={$f'(x)$},
  curve={-\x*\x+2*\x+3},dmin=-2,dmax=4,
  extra={\draw[dashed, line width=0.6pt] (-1,-6) -- (-1,5) (3,-6) -- (3,5);}]
\\[2mm]
% a possible f(x) = -x^3/3 + x^2 + 3x: minimum (-1, -5/3), maximum (3, 9).  Curve only with \showfp.
\sketch[xmin=-2,xmax=4,ymin=-4,ymax=10,unit=\figunit,yunit=0.35\figunit,xstep=\figxstep,ystep=\figystep,name={$f(x)$},
  curve={-\x*\x*\x/3+\x*\x+3*\x},show=\ifdefined\showfp 1\else 0\fi,dmin=-2,dmax=4,
  extra={\draw[dashed, line width=0.6pt] (-1,-4) -- (-1,10) (3,-4) -- (3,10);}]

\end{minipage}

\section{Versions that get no credit}\label{nocredit}
\begin{tabular}{@{}p{0.40\linewidth}p{0.54\linewidth}@{}}
\toprule
sketch & why it fails\\
\midrule
$f'$ with the same humps as $f$ & \tturns of $f$ must become zeros of $f'$\\
$f'$ of a parabola drawn curved & a parabola's gradient changes at a constant rate: a line\\
$f$ drawn through the zeros of $f'$ & zeros of $f'$ are the \tturns of $f$\\
key points not lined up & the marker checks the $x$-values against the given graph\\
\bottomrule
\end{tabular}

\section{What each grade looks like here}\label{grades}
\begin{tabular}{@{}p{0.14\linewidth}p{0.80\linewidth}@{}}
\toprule
Achieved & a correct sketch of $f'$ for a line or a parabola, zeros in the right places.\\
Merit & a correct $f'$ for a cubic or quartic (zeros, signs and shape), or a correct $f$ from $f'$.\\
Excellence & a full explanation linking the two graphs, e.g.\ why a sharp corner has no gradient, or why the maximum of $f'$ is not a maximum of $f$.\\
\bottomrule
\end{tabular}

\section{Checking with the fx-9750GIII}\label{calc}
\textsf{SET UP} (SHIFT MENU) $\to$ \textsf{Derivative: On}. In \textsf{GRAPH}, \textsf{TRACE} then shows the gradient $\dydx$ at each point.
Trace along $f$: the sign of the gradient tells you which side of the axis $f'$ is on.

\section{Ways to lose marks}\label{lose}
From the material plus teaching observation.
\begin{enumerate}[leftmargin=7mm, itemsep=2pt]
\item[\textbf{\lmnum{tpcopy}}] \lmtxt{tpcopy}
\item[\textbf{\lmnum{sign}}] \lmtxt{sign}
\item[\textbf{\lmnum{degree}}] \lmtxt{degree}
\item[\textbf{\lmnum{steep}}] \lmtxt{steep}
\item[\textbf{\lmnum{roots}}] \lmtxt{roots}
\item[\textbf{\lmnum{marks}}] \lmtxt{marks}
\end{enumerate}
\end{document}
